\documentclass[12pt,a4paper]{article}
\usepackage[utf8]{inputenc}
\usepackage[spanish]{babel}
\usepackage{geometry}
\usepackage{enumitem}
\usepackage{xcolor}
\usepackage{tcolorbox}
\usepackage{hyperref}
\usepackage{fancyhdr}
\usepackage{graphicx}
\usepackage{amsmath}

\geometry{margin=2.5cm}
\pagestyle{fancy}
\fancyhf{}
\rhead{NLPY-1}
\lhead{Proyecto Final Libre}
\rfoot{\thepage}

\definecolor{nablue}{RGB}{0,102,204}
\definecolor{nagreen}{RGB}{34,139,34}
\definecolor{naorange}{RGB}{255,140,0}

\title{\textbf{\Huge Proyecto Final Libre}\\[0.5em]
\Large Curso de Programación en Python NLPY-1\\[0.5em]
\large Semana 9}
\author{nablabs}
\date{2025}

\begin{document}

\maketitle
\thispagestyle{empty}

\vspace{1cm}

\begin{tcolorbox}[colback=nablue!5!white,colframe=nablue,title=\textbf{Objetivo del Proyecto}]
Desarrollar una aplicación completa en Python que integre \textbf{todos} los conceptos y técnicas aprendidas durante el curso, demostrando dominio de la programación y capacidad de resolver problemas reales.
\end{tcolorbox}

\newpage

\tableofcontents

\newpage

\section{Descripción General}

Este proyecto final representa la culminación de tu aprendizaje en el curso NLPY-1. A diferencia de los proyectos semanales, este proyecto es \textbf{libre}: tú decides qué aplicación desarrollar, siempre y cuando cumpla con los requisitos técnicos establecidos.

\subsection{Libertad Creativa}

Tienes total libertad para elegir:
\begin{itemize}
    \item El tema de tu aplicación
    \item El dominio del problema a resolver
    \item La interfaz de usuario (consola, GUI, web)
    \item Las librerías adicionales que desees utilizar
    \item El nivel de complejidad y características extras
\end{itemize}

\subsection{Salón de la Fama}

\begin{tcolorbox}[colback=nagreen!5!white,colframe=nagreen,title=\textbf{¡Oportunidad Especial!}]
Los proyectos más \textbf{sorprendentes} y \textbf{llamativos} serán presentados en el \textbf{Salón de la Fama} de la página web de nablabs. Tu trabajo podría inspirar a futuros estudiantes y formar parte de la historia del curso.

\vspace{0.5em}
\textit{Criterios para el Salón de la Fama:}
\begin{itemize}
    \item Originalidad y creatividad
    \item Calidad técnica de implementación
    \item Impacto visual o funcional
    \item Utilidad práctica
    \item Código limpio y bien documentado
\end{itemize}
\end{tcolorbox}

\section{Requisitos Técnicos Obligatorios}

Tu proyecto \textbf{DEBE} implementar todos los temas vistos durante el curso:

\subsection{Semana 0: Funciones y Variables}
\begin{itemize}
    \item Uso de funciones con parámetros y valores de retorno
    \item Variables de diferentes tipos (strings, integers, floats)
    \item Formateo de strings (f-strings)
    \item Métodos de strings (strip, title, capitalize, split, etc.)
    \item Operadores aritméticos
    \item Conversión de tipos (type casting)
    \item Función \texttt{input()} para interacción con el usuario
\end{itemize}

\subsection{Semana 1: Condicionales}
\begin{itemize}
    \item Sentencias \texttt{if}, \texttt{elif}, \texttt{else}
    \item Operadores de comparación (\texttt{>}, \texttt{<}, \texttt{>=}, \texttt{<=}, \texttt{==}, \texttt{!=})
    \item Operadores lógicos (\texttt{and}, \texttt{or}, \texttt{not})
    \item Valores booleanos y expresiones booleanas
    \item (Opcional) Operador \texttt{match} para casos complejos
\end{itemize}

\subsection{Semana 2: Bucles}
\begin{itemize}
    \item Bucles \texttt{while} para iteraciones condicionales
    \item Bucles \texttt{for} para iteración sobre secuencias
    \item Uso de \texttt{range()}
    \item Sentencias \texttt{break} y \texttt{continue}
    \item Validación de entrada con bucles
\end{itemize}

\subsection{Semana 3: Excepciones}
\begin{itemize}
    \item Bloques \texttt{try} y \texttt{except}
    \item Manejo de múltiples tipos de excepciones
    \item Uso de \texttt{else} y \texttt{finally}
    \item Creación de excepciones personalizadas con \texttt{raise}
    \item Mensajes de error descriptivos
\end{itemize}

\subsection{Semana 4: Librerías}
\begin{itemize}
    \item Importación de módulos estándar (\texttt{random}, \texttt{datetime}, \texttt{statistics}, etc.)
    \item Uso de al menos 2 librerías externas instaladas con \texttt{pip}
    \item Creación de al menos 1 módulo propio
    \item Uso apropiado de \texttt{import} y \texttt{from ... import ...}
\end{itemize}

\subsection{Semana 5: Pruebas Unitarias}
\begin{itemize}
    \item Archivo separado de pruebas usando \texttt{pytest}
    \item Al menos 10 pruebas unitarias diferentes
    \item Pruebas para casos normales, casos límite y excepciones
    \item Uso de \texttt{assert} para validaciones
    \item Todas las pruebas deben pasar exitosamente
\end{itemize}

\subsection{Semana 6: Entrada/Salida de Archivos}
\begin{itemize}
    \item Lectura y escritura de archivos
    \item Uso del bloque \texttt{with} para manejo seguro de archivos
    \item Trabajo con archivos CSV (usando el módulo \texttt{csv})
    \item Persistencia de datos entre ejecuciones
    \item Manejo de errores en operaciones de archivos
\end{itemize}

\subsection{Semana 7: Expresiones Regulares}
\begin{itemize}
    \item Uso del módulo \texttt{re}
    \item Al menos 3 patrones regex diferentes
    \item Validación de formatos (emails, teléfonos, cédulas, etc.)
    \item Uso de funciones como \texttt{search()}, \texttt{match()}, \texttt{findall()}, o \texttt{sub()}
    \item Patrones con metacaracteres y cuantificadores
\end{itemize}

\subsection{Semana 8: Programación Orientada a Objetos}
\begin{itemize}
    \item Al menos 4 clases diferentes bien diseñadas
    \item Uso de \texttt{\_\_init\_\_} (constructor)
    \item Atributos de instancia y métodos
    \item Al menos 1 ejemplo de herencia
    \item Uso de \texttt{@property} para encapsulación
    \item Métodos especiales (\texttt{\_\_str\_\_}, \texttt{\_\_repr\_\_})
    \item (Opcional) \texttt{@classmethod} o \texttt{@staticmethod}
\end{itemize}

\subsection{Semana 9: Et Cetera}
\begin{itemize}
    \item Uso de conjuntos (sets) o diccionarios avanzados
    \item Al menos 2 list comprehensions
    \item Al menos 1 dictionary comprehension
    \item Uso de funciones lambda
    \item (Opcional) Uso de \texttt{*args} o \texttt{**kwargs}
    \item (Opcional) Decoradores o generadores
\end{itemize}

\section{Evaluación}

\subsection{Distribución de Puntos}

\begin{center}
\begin{tabular}{|l|c|}
\hline
\textbf{Criterio} & \textbf{Puntos} \\
\hline
Implementación de todos los requisitos técnicos & 40 \\
Funcionalidad y lógica del programa & 20 \\
Calidad del código (limpieza, organización) & 15 \\
Documentación y comentarios & 10 \\
Pruebas unitarias completas y efectivas & 10 \\
Creatividad e innovación & 5 \\
\hline
\textbf{Total} & \textbf{100} \\
\hline
\textbf{Bonus: Juego Retro} & \textbf{+20} \\
\hline
\end{tabular}
\end{center}

\subsection{Bonus: Juego Retro (+20 puntos)}

\begin{tcolorbox}[colback=naorange!5!white,colframe=naorange,title=\textbf{Desafío Especial: Juegos Retro}]
Si decides desarrollar un \textbf{juego retro clásico}, recibirás \textbf{20 puntos adicionales} por la complejidad técnica que esto implica.

\vspace{0.5em}
\textbf{Ejemplos de juegos retro:}
\begin{itemize}
    \item Space Invaders
    \item Bomberman (estilo NES)
    \item Donkey Kong
    \item Snake
    \item Pac-Man
    \item Asteroids
    \item Pong
    \item Tetris
\end{itemize}

\vspace{0.5em}
\textbf{Requisitos adicionales para el bonus:}
\begin{itemize}
    \item El juego debe ser jugable y funcional
    \item Debe incluir mecánicas básicas del juego original
    \item Interfaz gráfica (puede usar \texttt{pygame}, \texttt{tkinter}, etc.)
    \item Sistema de puntuación
    \item \textbf{IMPORTANTE:} Debe implementar TODOS los temas del curso
\end{itemize}

\vspace{0.5em}
\textit{Nota: Desarrollar un juego retro es significativamente más complejo, pero también más gratificante y tiene mayor probabilidad de llegar al Salón de la Fama.}
\end{tcolorbox}

\section{Ideas de Proyectos Sugeridas}

A continuación se presentan 8 ideas de proyectos completamente desarrolladas. Puedes elegir una de estas o crear tu propia idea original.

\subsection{Sistema de Gestión de Citas Médicas}

\textbf{Descripción:} Aplicación para gestionar citas médicas en una clínica u hospital.

\vspace{0.5em}
\textbf{Componentes principales:}
\begin{itemize}
    \item \textbf{Clases:} Paciente, Doctor, Cita, Clínica
    \item \textbf{GUI:} Formularios para registrar pacientes/doctores, calendario visual de citas
    \item \textbf{Archivos:} Guardar datos en CSV/JSON, exportar reportes de citas
    \item \textbf{Regex:} Validar cédulas costarricenses, emails, teléfonos
    \item \textbf{Librerías:} \texttt{datetime} para manejo de fechas y horarios, \texttt{random} para generar códigos únicos de cita
    \item \textbf{Testing:} Validar lógica de reservas, detectar conflictos de horarios
\end{itemize}

\vspace{0.5em}
\textbf{Funcionalidades:}
\begin{itemize}
    \item Registro de pacientes y doctores
    \item Agendar, modificar y cancelar citas
    \item Validar disponibilidad de horarios
    \item Búsqueda de citas por paciente, doctor o fecha
    \item Recordatorios de citas próximas
    \item Generación de reportes estadísticos
\end{itemize}

\subsection{Gestor de Finanzas Personales}

\textbf{Descripción:} Aplicación para controlar ingresos, gastos y presupuestos personales.

\vspace{0.5em}
\textbf{Componentes principales:}
\begin{itemize}
    \item \textbf{Clases:} Cuenta, Transacción, Categoría, Usuario
    \item \textbf{GUI:} Dashboard con gráficos de gastos, formularios de ingreso/egreso
    \item \textbf{Archivos:} Persistencia de transacciones, importar extractos bancarios CSV
    \item \textbf{Regex:} Validar números de cuenta, parsear descripciones de transacciones
    \item \textbf{Librerías:} \texttt{matplotlib} para gráficos, \texttt{statistics} para análisis
    \item \textbf{Excepciones:} Manejo de saldos insuficientes, transacciones inválidas
\end{itemize}

\vspace{0.5em}
\textbf{Funcionalidades:}
\begin{itemize}
    \item Registro de ingresos y gastos
    \item Categorización de transacciones
    \item Visualización de gastos por categoría
    \item Presupuestos mensuales con alertas
    \item Reportes financieros con gráficos
    \item Importación de datos desde archivos CSV
\end{itemize}

\subsection{Sistema de Biblioteca Digital}

\textbf{Descripción:} Plataforma para gestionar préstamos de libros en una biblioteca.

\vspace{0.5em}
\textbf{Componentes principales:}
\begin{itemize}
    \item \textbf{Clases:} Libro, Usuario, Préstamo, Biblioteca
    \item \textbf{GUI:} Catálogo visual, búsqueda avanzada, gestión de préstamos
    \item \textbf{Archivos:} Base de datos de libros en CSV, historial de préstamos
    \item \textbf{Regex:} Búsqueda por patrones (ISBN, autor, título)
    \item \textbf{Librerías:} \texttt{datetime} para fechas de devolución, \texttt{random} para códigos
    \item \textbf{Testing:} Validar reglas de préstamo, cálculo de multas por atraso
\end{itemize}

\vspace{0.5em}
\textbf{Funcionalidades:}
\begin{itemize}
    \item Catálogo de libros con búsqueda avanzada
    \item Sistema de préstamos y devoluciones
    \item Control de disponibilidad
    \item Cálculo automático de multas
    \item Historial de préstamos por usuario
    \item Reservas de libros prestados
\end{itemize}

\subsection{Administrador de Tareas y Proyectos}

\textbf{Descripción:} Aplicación tipo Trello para gestionar tareas y proyectos.

\vspace{0.5em}
\textbf{Componentes principales:}
\begin{itemize}
    \item \textbf{Clases:} Tarea, Proyecto, Etiqueta, Usuario
    \item \textbf{GUI:} Tableros visuales, filtros, calendario de tareas
    \item \textbf{Archivos:} Guardar proyectos y tareas, exportar reportes
    \item \textbf{Regex:} Parsear fechas, buscar por etiquetas
    \item \textbf{Librerías:} \texttt{datetime}, \texttt{os} para organizar archivos
    \item \textbf{Comprehensions:} Filtrar tareas por estado, prioridad, fecha
\end{itemize}

\vspace{0.5em}
\textbf{Funcionalidades:}
\begin{itemize}
    \item Creación de proyectos y tareas
    \item Estados de tareas (por hacer, en progreso, completado)
    \item Sistema de prioridades
    \item Fechas límite con recordatorios
    \item Etiquetas y categorías personalizadas
    \item Vista de calendario
    \item Exportación de reportes de productividad
\end{itemize}

\subsection{Sistema de Reservas de Hotel/Hostel}

\textbf{Descripción:} Plataforma para gestionar reservas de habitaciones.

\vspace{0.5em}
\textbf{Componentes principales:}
\begin{itemize}
    \item \textbf{Clases:} Habitación, Reserva, Cliente, Hotel
    \item \textbf{GUI:} Calendario de disponibilidad, formulario de reserva
    \item \textbf{Archivos:} Persistencia de reservas, generación de facturas
    \item \textbf{Regex:} Validar pasaportes, tarjetas de crédito (formato)
    \item \textbf{Librerías:} \texttt{datetime} para fechas, \texttt{random} para códigos de reserva
    \item \textbf{Excepciones:} Manejo de overbooking, cancelaciones
\end{itemize}

\vspace{0.5em}
\textbf{Funcionalidades:}
\begin{itemize}
    \item Búsqueda de habitaciones disponibles
    \item Sistema de reservas con confirmación
    \item Diferentes tipos de habitaciones y tarifas
    \item Check-in y check-out
    \item Generación de facturas
    \item Cancelaciones y reembolsos
    \item Estadísticas de ocupación
\end{itemize}

\subsection{Plataforma de Quiz/Exámenes Online}

\textbf{Descripción:} Sistema para crear y realizar exámenes en línea.

\vspace{0.5em}
\textbf{Componentes principales:}
\begin{itemize}
    \item \textbf{Clases:} Pregunta, Examen, Estudiante, Resultado
    \item \textbf{GUI:} Interfaz de examen, timer, resultados con gráficos
    \item \textbf{Archivos:} Banco de preguntas en CSV, exportar calificaciones
    \item \textbf{Regex:} Validar formatos de respuestas específicas
    \item \textbf{Librerías:} \texttt{random} para barajear preguntas, \texttt{statistics} para análisis
    \item \textbf{Testing:} Verificar lógica de calificación automática
\end{itemize}

\vspace{0.5em}
\textbf{Funcionalidades:}
\begin{itemize}
    \item Creación de bancos de preguntas
    \item Diferentes tipos de preguntas (opción múltiple, verdadero/falso, respuesta corta)
    \item Temporizador para exámenes
    \item Calificación automática
    \item Barajeo de preguntas y opciones
    \item Estadísticas de desempeño
    \item Exportación de resultados
\end{itemize}

\subsection{Sistema de Inventario para Pulpería/Minisuper}

\textbf{Descripción:} Aplicación de punto de venta e inventario para negocio.

\vspace{0.5em}
\textbf{Componentes principales:}
\begin{itemize}
    \item \textbf{Clases:} Producto, Venta, Cliente, Inventario
    \item \textbf{GUI:} Punto de venta, gestión de stock, reportes visuales
    \item \textbf{Archivos:} Base de datos de productos, registro de ventas
    \item \textbf{Regex:} Validar códigos de barras, buscar productos
    \item \textbf{Librerías:} \texttt{datetime} para ventas, \texttt{statistics} para análisis
    \item \textbf{Excepciones:} Stock insuficiente, precios inválidos
\end{itemize}

\vspace{0.5em}
\textbf{Funcionalidades:}
\begin{itemize}
    \item Registro de productos con código de barras
    \item Control de inventario en tiempo real
    \item Proceso de venta con cálculo de total
    \item Alertas de stock bajo
    \item Historial de ventas
    \item Reportes de productos más vendidos
    \item Gestión de proveedores
\end{itemize}

\subsection{Aplicación de Recetas de Cocina Tica}

\textbf{Descripción:} Plataforma para gestionar y buscar recetas de comida costarricense.

\vspace{0.5em}
\textbf{Componentes principales:}
\begin{itemize}
    \item \textbf{Clases:} Receta, Ingrediente, Usuario, Favoritos
    \item \textbf{GUI:} Buscador visual, filtros por ingredientes, timer de cocina
    \item \textbf{Archivos:} Base de datos de recetas, exportar lista de compras
    \item \textbf{Regex:} Parsear cantidades (``2 tazas'', ``500g''), buscar ingredientes
    \item \textbf{Librerías:} \texttt{random} para sugerir recetas del día
    \item \textbf{Comprehensions:} Filtrar recetas por ingredientes disponibles
\end{itemize}

\vspace{0.5em}
\textbf{Funcionalidades:}
\begin{itemize}
    \item Base de datos de recetas típicas costarricenses
    \item Búsqueda por nombre, ingredientes o categoría
    \item Sistema de favoritos
    \item Generación automática de lista de compras
    \item Conversión de unidades
    \item Timer para cocinar
    \item Sugerencias aleatorias de recetas
    \item Exportar recetas a archivo
\end{itemize}

\section{Entregables}

\subsection{Código Fuente}

\begin{itemize}
    \item Todos los archivos \texttt{.py} necesarios
    \item Módulos organizados en carpetas apropiadas
    \item Archivo de pruebas \texttt{test\_*.py}
    \item Archivos de datos (CSV, JSON, etc.)
    \item Archivo \texttt{requirements.txt} con todas las dependencias
\end{itemize}

\subsection{Documentación}

\begin{itemize}
    \item \texttt{README.md} con:
    \begin{itemize}
        \item Descripción del proyecto
        \item Instrucciones de instalación
        \item Cómo ejecutar el programa
        \item Cómo ejecutar las pruebas
        \item Ejemplos de uso
    \end{itemize}
    \item Comentarios en el código explicando lógica compleja
    \item Docstrings en todas las funciones y clases
\end{itemize}

\subsection{Video Demostración (Opcional pero Recomendado)}

\begin{itemize}
    \item Video de 3-5 minutos mostrando la funcionalidad
    \item Explicación de características principales
    \item Demostración de casos de uso
    \item \textit{Aumenta significativamente las posibilidades de llegar al Salón de la Fama}
\end{itemize}

\section{Cronograma de Entrega}

\begin{center}
\begin{tabular}{|l|l|}
\hline
\textbf{Fecha} & \textbf{Entregable} \\
\hline
Semana 9, Día 3 & Propuesta de proyecto (1 párrafo) \\
Semana 9, Día 5 & Diseño de clases y estructura \\
Semana 10, Día 3 & Avance del 50\% funcional \\
Semana 10, Día 7 & \textbf{Entrega final completa} \\
\hline
\end{tabular}
\end{center}

\section{Recomendaciones}

\begin{enumerate}
    \item \textbf{Planifica primero:} Diseña tu aplicación antes de programar. Define las clases, funciones y flujo de datos.
    
    \item \textbf{Desarrollo iterativo:} No intentes hacer todo de una vez. Desarrolla por características incrementales.
    
    \item \textbf{Prueba continuamente:} Escribe pruebas mientras desarrollas, no al final.
    
    \item \textbf{Commits frecuentes:} Si usas Git, haz commits pequeños y frecuentes con mensajes descriptivos.
    
    \item \textbf{Código limpio:} Sigue convenciones de estilo (PEP 8), usa nombres descriptivos, evita repetición.
    
    \item \textbf{Maneja errores:} No dejes que tu programa crashee. Maneja todas las excepciones apropiadamente.
    
    \item \textbf{Documenta:} Escribe docstrings y comentarios útiles. Tu código debe ser entendible por otros.
    
    \item \textbf{Sé creativo:} Los proyectos más memorables son los que muestran originalidad e innovación.
    
    \item \textbf{Piensa en el usuario:} Haz que tu aplicación sea intuitiva y fácil de usar.
    
    \item \textbf{Ve más allá:} Los requisitos son el mínimo. Agregar características extra te distinguirá.
\end{enumerate}

\section{Criterios de Calidad del Código}

\subsection{Organización}
\begin{itemize}
    \item Estructura de archivos lógica y clara
    \item Separación de responsabilidades (clases bien definidas)
    \item Módulos importados correctamente
\end{itemize}

\subsection{Estilo}
\begin{itemize}
    \item Cumplimiento con PEP 8
    \item Nombres de variables descriptivos (no \texttt{x}, \texttt{y}, \texttt{data})
    \item Funciones con un propósito claro y único
    \item Longitud apropiada de funciones (no más de 50 líneas idealmente)
\end{itemize}

\subsection{Documentación}
\begin{itemize}
    \item Docstrings en todas las clases y funciones públicas
    \item Comentarios explicando lógica no obvia
    \item README completo y claro
\end{itemize}

\subsection{Manejo de Errores}
\begin{itemize}
    \item Excepciones manejadas apropiadamente
    \item Mensajes de error útiles para el usuario
    \item No uso de \texttt{except:} sin especificar el tipo de excepción
\end{itemize}

\section{Ejemplos de Implementación de Requisitos}

\subsection{Ejemplo: Expresiones Regulares}

\begin{verbatim}
import re

class ValidadorCostarricense:
    """Validador de formatos costarricenses"""
    
    @staticmethod
    def validar_cedula(cedula):
        """Valida formato de cédula costarricense (0-0000-0000)"""
        patron = r'^\d-\d{4}-\d{4}$'
        return re.match(patron, cedula) is not None
    
    @staticmethod
    def validar_telefono(telefono):
        """Valida teléfono costarricense (8 dígitos)"""
        patron = r'^[2-8]\d{3}-?\d{4}$'
        return re.match(patron, telefono) is not None
\end{verbatim}

\subsection{Ejemplo: List Comprehension}

\begin{verbatim}
# Filtrar tareas pendientes con alta prioridad
tareas_urgentes = [tarea for tarea in self.tareas 
                   if tarea.prioridad == 'alta' 
                   and tarea.estado == 'pendiente']

# Obtener nombres de clientes activos
nombres_activos = [cliente.nombre.title() 
                   for cliente in self.clientes 
                   if cliente.activo]
\end{verbatim}

\subsection{Ejemplo: Testing}

\begin{verbatim}
import pytest
from sistema import Cuenta, TransaccionInvalida

def test_deposito_valido():
    cuenta = Cuenta("Juan", 1000)
    cuenta.depositar(500)
    assert cuenta.saldo == 1500

def test_retiro_saldo_insuficiente():
    cuenta = Cuenta("María", 100)
    with pytest.raises(TransaccionInvalida):
        cuenta.retirar(200)

def test_transferencia_exitosa():
    cuenta1 = Cuenta("Pedro", 1000)
    cuenta2 = Cuenta("Ana", 500)
    cuenta1.transferir(cuenta2, 300)
    assert cuenta1.saldo == 700
    assert cuenta2.saldo == 800
\end{verbatim}

\section{Preguntas Frecuentes}

\subsection{¿Puedo usar librerías que no vimos en clase?}

Sí, de hecho es recomendado. Debes usar al menos 2 librerías externas, y puedes elegir las que mejor se adapten a tu proyecto.

\subsection{¿Tengo que usar GUI o puede ser solo consola?}

Puedes usar consola, GUI o incluso web. Lo importante es la funcionalidad y el cumplimiento de requisitos técnicos.

\subsection{¿Puedo trabajar en equipo?}

Este proyecto es individual. Sin embargo, puedes discutir ideas y conceptos con compañeros, siempre que el código sea tuyo.

\subsection{¿Qué pasa si mi proyecto no incluye algún requisito?}

Cada requisito faltante resultará en deducción de puntos. Asegúrate de implementar TODO.

\subsection{¿Puedo hacer algo diferente a las ideas sugeridas?}

¡Por supuesto! Las ideas son solo sugerencias. Tu creatividad es bienvenida y valorada.

\subsection{¿Cuánto código debo escribir?}

No hay un mínimo estricto, pero un proyecto completo típicamente tiene 500-1500 líneas de código (sin contar archivos de datos).

\section{Recursos Adicionales}

\subsection{Documentación Oficial}
\begin{itemize}
    \item Python: \url{https://docs.python.org/3/}
    \item Pytest: \url{https://docs.pytest.org/}
    \item Expresiones Regulares: \url{https://docs.python.org/3/library/re.html}
\end{itemize}

\subsection{Librerías Útiles}
\begin{itemize}
    \item \texttt{tkinter} - GUI básica
    \item \texttt{pygame} - Juegos 2D
    \item \texttt{matplotlib} - Gráficos
    \item \texttt{pandas} - Manipulación de datos
    \item \texttt{requests} - APIs HTTP
    \item \texttt{pillow} - Procesamiento de imágenes
\end{itemize}

\section{Checklist Final}

Antes de entregar, verifica que tu proyecto incluya:

\begin{itemize}
    \item[$\square$] Funciones con parámetros y return
    \item[$\square$] Variables y type casting
    \item[$\square$] F-strings y métodos de strings
    \item[$\square$] Condicionales (if/elif/else)
    \item[$\square$] Operadores lógicos
    \item[$\square$] Bucles while y for
    \item[$\square$] Break y continue
    \item[$\square$] Try/except con al menos 3 tipos de excepciones
    \item[$\square$] Raise para excepciones personalizadas
    \item[$\square$] Importación de módulos estándar
    \item[$\square$] Al menos 2 librerías externas
    \item[$\square$] Creación de módulo propio
    \item[$\square$] Archivo de pruebas con pytest
    \item[$\square$] Al menos 10 pruebas unitarias
    \item[$\square$] Lectura y escritura de archivos
    \item[$\square$] Uso de archivos CSV
    \item[$\square$] Bloque with para archivos
    \item[$\square$] Al menos 3 patrones regex
    \item[$\square$] Funciones del módulo re
    \item[$\square$] Al menos 4 clases
    \item[$\square$] Uso de \_\_init\_\_
    \item[$\square$] Al menos 1 ejemplo de herencia
    \item[$\square$] @property
    \item[$\square$] \_\_str\_\_ y/o \_\_repr\_\_
    \item[$\square$] Al menos 2 list comprehensions
    \item[$\square$] Al menos 1 dictionary comprehension
    \item[$\square$] Funciones lambda
    \item[$\square$] Sets o diccionarios avanzados
    \item[$\square$] README.md completo
    \item[$\square$] requirements.txt
    \item[$\square$] Código comentado y documentado
    \item[$\square$] Todas las pruebas pasan
\end{itemize}

\vspace{1cm}

\begin{tcolorbox}[colback=nablue!5!white,colframe=nablue,title=\textbf{¡Mucho Éxito!}]
Este proyecto es tu oportunidad para demostrar todo lo que has aprendido. No tengas miedo de ser ambicioso y creativo. Los mejores proyectos son aquellos que resuelven problemas reales o que simplemente son divertidos de usar.

\vspace{0.5em}
\textbf{Recuerda:} El objetivo no es solo cumplir requisitos, sino crear algo de lo que estés orgulloso.

\vspace{0.5em}
\centering
\textit{¡Nos vemos en el Salón de la Fama!}
\end{tcolorbox}

\end{document}
